%!TEX program = lualatex

\documentclass[11pt]{article}

\usepackage[fleqn,tbtags]{amsmath}
\usepackage[ukrainian]{babel}
\usepackage{booktabs}
\usepackage{fontspec}
\usepackage{geometry}
\usepackage{graphicx}
\usepackage{hyperref}
\usepackage{mathtools}
\usepackage{microtype}
\usepackage[cache=false, outputdir=../out]{minted}
\usepackage{multicol}
\usepackage{multirow}
\usepackage[parfill]{parskip}
\usepackage[sfdefault]{roboto}

\newcommand{\TITLE}{Лабораторна робота №2}
\newcommand{\DISCIPLINE}{Теорія біомедичних сигналів}
\newcommand{\SUBJECT}{Візуалізація сигналів з використанням Python}
\newcommand{\STUDENT}{}
\newcommand{\TEACHER}{}

\emergencystretch 5em
\geometry{
    a4paper,
    left=1.5cm,
    right=1.5cm,
    top=1.5cm,
    bottom=2cm,
}

\setmonofont{Roboto}
\setminted{fontsize=\footnotesize}

\newtheorem{theorem}{Теорема}

\begin{document}
\clearpage\begin{titlepage}
    \begin{minipage}{\textwidth}
        \begin{minipage}{0.3\textwidth}
            \includegraphics[width=\textwidth]{./images/kpi}
        \end{minipage}
        \hfill
        \begin{minipage}{0.7\textwidth}
            \centering
            \Large
            МІНІСТЕРСТВО ОСВІТИ І НАУКИ УКРАЇНИ \\
            НАЦІОНАЛЬНИЙ ТЕХНІЧНИЙ УНІВЕРСИТЕТ УКРАЇНИ \\
            «КИЇВСЬКИЙ ПОЛІТЕХНІЧНИЙ ІНСТИТУТ \\
            ІМЕНІ ІГОРЯ СІКОРСЬКОГО» \\
            ФАКУЛЬТЕТ БІОМЕДИЧНОЇ ІНЖЕНЕРІЇ \\
            КАФЕДРА БІОМЕДИЧНОЇ КІБЕРНЕТИКИ
        \end{minipage}
    \end{minipage}

    \vfill

    \begin{center}
        \LARGE
        \textbf{\TITLE} \\
        з дисципліни \textbf{«\DISCIPLINE»} \\
        на тему: «\SUBJECT»
    \end{center}

    \vfill

    \begin{minipage}{\textwidth}
        \hfill
        \begin{minipage}{0.4\textwidth}
            \Large

            \textbf{Виконав:} \\
            \STUDENT \\
            \textbf{Перевірив:} \\
            \TEACHER \\

            Зараховано від:
            \hrulefill \\
            \begin{center}
                \small
                \hrulefill \\
                (підпис викладача)
            \end{center}
        \end{minipage}
    \end{minipage}

    \vfill

    \begin{center}
        \LARGE
        Київ-\the\year{}
    \end{center}
\end{titlepage}

\clearpage\tableofcontents
\clearpage\begin{abstract}
    \section*{Мета роботи}
    \label{sec:meta}

    \begin{enumerate}
        \large
        \item навчитись застосовувати метод Сааті в практичних задачах
    \end{enumerate}

    \section*{Завдання}
    \label{sec:tasks}

    \begin{enumerate}
        \large
        \item Навчитись визначати найкращу альтернативу за деякою кількістю критеріїв
        \item Визначити наближені значення елементів власних векторів
        \item Описати пріоритети критеріїв та альтернатив
    \end{enumerate}
\end{abstract}

\clearpage\section{Частина 1}
\label{sec:task1}

\subsection{Завдання}
\label{subsec:task1_task}

Вимоги до програми:
\begin{itemize}
      \item Забезпечує багаторазову генерацію масивів
            \begin{align}
                  X_M = \{x_n^0, n = 1, \dots, N\}, M = 1, \dots, M_0
            \end{align}
            незалежних однаково випадкових величин $x_n^0$ рівномірно розподілених
            в інтервалі значень $x_{min}^0, x_{max}^0$, що імітують значення критерію,
            який характеризує альтернативи $d \in \{D_1, \dots, D_N\}$
      \item Визначає абсолютного лідера окремого експерименту за максимальним
            значенням критерію $x^* = max x_n^0, 1 \leq n \leq N$
      \item Визначає умовного лідера за максимальним значенням критерію в інтервалі
            індексів $[1, t], t < N$, тобто визначає число
            \begin{align}
                  x_{1, t}^{max} = max \{x_1^0, \dots, x_t^0\}
            \end{align}
      \item Остаточним рішення вважається перша з альтернатив $\widetilde{d^*}$
            в діапазоні індексів $t, N$, яку характеризує критерій,
            що перевищує число $x_{1, t}^{max}$
      \item Проводиться $M_o = 1000$ окремих експериментів (цикл), для кожного з
            котрих оцінюють ймовірність правильного рішення для кожного
            $t \in [1, N]$ (внутрішній цикл), тобто ймовірність вибору альтернативи
            $\widetilde{d^*} \in D$, яка задовольняє умову
            $|x^0(\widetilde{d^*}) - x^0(d^*)| \leq \Delta$
      \item Ймовірність $\widetilde{P}$ оцінюють за формулою
            \begin{align}
                  P(t) = \frac{m(t)}{M_0}
            \end{align}
            де $m(t)$ - число успіхів на $t$-му кроці в серії з $M_0$ експериментів
      \item За максимумом обчисленої ймовірністі $P(t)$ визначають момент
            оптимальної зупинки $t^*$ при різних значеннях поступки $\Delta \geq 0$,
            яку задають у відсотках від абсолютного максимуму
            $x^* = max x_n^0, 1 \leq n \leq N$
      \item Будують графіки залежності:
            \begin{itemize}
                  \item ймовірністі $P(t)$ від номеру кроку $t$ за різних $\Delta \geq 0$
                  \item ймовірністі $\widetilde{P}$ від величини поступки $\Delta \geq 0$
                  \item номеру зупинку $t^*$ від величини поступки $\Delta \geq 0$
            \end{itemize}
\end{itemize}

\subsection{Код програми}
\label{subsec:task1_code}
\inputminted{python}{../src/task1.py}

\subsection{Результати}
\label{subsec:task1_results}

\begin{figure}[!ht]
      \centering
      \includegraphics[width=\textwidth]{./images/task1.png}
      \caption{Результати статистичного експерименту}
      \label{fig:graphs}
\end{figure}

На графіку бачимо як за збільшення значення поступки $\Delta$
оптимальний момент $t^* \in [1, \dots, N]$, коли слід приймати остаточне рішення
зменшується, а вірогідність $\widetilde{P}$ прийняти правильне рішення збільшується,
що й треба було доказати.

\begin{tabular}{|c|c|}
      \toprule
      $\Delta$ & $P$   \\

      \midrule
      0        & 0.409 \\
      \hline
      0.02     & 0.675 \\
      \hline
      0.04     & 0.796 \\
      \hline
      0.06     & 0.856 \\
      \hline
      0.08     & 0.898 \\
      \hline
      0.1      & 0.915 \\

      \bottomrule
\end{tabular}

\clearpage\section{Завдання 2}
\label{sec:task2}

\subsection{Код Программи}
\label{subsec:task2_code}
\inputminted{python}{../src/task2.py}

\subsection{Результати}
\label{subsec:task2_results}

\clearpage\section{Частина 3}
\label{sec:task3}

\subsection{Завдання}
\label{subsec:task3_task}

Програмна реалізація статистичного експерименту,
що забезпечує оцінювання властивостей модифікованого
методу оптимальної зупинки.

\subsection{Код програми}
\label{subsec:task3_code}
\inputminted{python}{../src/task1.py}

\subsection{Результат}
\label{subsec:task3_result}

\begin{figure}[!ht]
    \centering
    \includegraphics[width=\textwidth]{./images/task1.png}
    \caption{Результати статистичного експерименту}
    \label{fig:graphs}
\end{figure}

\begin{tabular}{|c|c|}
    \toprule
    $\Delta$ & $P$   \\

    \midrule
    0        & 0.409 \\
    \hline
    0.02     & 0.675 \\
    \hline
    0.04     & 0.796 \\
    \hline
    0.06     & 0.856 \\
    \hline
    0.08     & 0.898 \\
    \hline
    0.1      & 0.915 \\

    \bottomrule
\end{tabular}

\clearpage\section{Частина 4}
\label{sec:task4}

\subsection{Завдання}
\label{subsec:task4_task}

Для сигналів з п. 2 лабораторної роботи №4 побудувати спектрограми
\begin{enumerate}
    \item з вікном тривалості 0,1 с та 2 с (без перекриття);
    \item з вікном тривалості 1 с з перекриттям 50\%.
\end{enumerate}
Застосувати функції colorbar для візуалізації значень спектрограми.
Зробити висновки щодо відображення часових властивостей сигналів у
спектрограмі. Порівняти інформативність спектрограм та спектрів з п. 2
лабораторної роботи №4. Побудувати тривимірні графіки двох різних
спектрограм з отриманих.

\subsection{Код программи}
\label{subsec:task4_code}
\inputminted{python}{../src/task4.py}

\subsection{Результати}
\label{subsec:task4_results}

Спектрограми дають дозволяють точно відстежити кількість, протяжність та частоту
складових сигналу, чого не можна сказати про амплітудний спектр.

\begin{figure}[!ht]
    \centering
    \includegraphics[width=\textwidth]{./images/task4_0.png}
    % \caption{}
    % \label{fig:}
\end{figure}

\begin{figure}[!ht]
    \centering
    \includegraphics[width=\textwidth]{./images/task4_1.png}
    % \caption{}
    % \label{fig:}
\end{figure}

\begin{figure}[!ht]
    \centering
    \includegraphics[width=\textwidth]{./images/task4_2.png}
    % \caption{}
    % \label{fig:}
\end{figure}

\begin{figure}[!ht]
    \centering
    \includegraphics[width=\textwidth]{./images/task4_3.png}
    % \caption{}
    % \label{fig:}
\end{figure}

\begin{figure}[!ht]
    \centering
    \includegraphics[width=\textwidth]{./images/task4_4.png}
    % \caption{}
    % \label{fig:}
\end{figure}

\begin{figure}[!ht]
    \centering
    \includegraphics[width=\textwidth]{./images/task4_5.png}
    % \caption{}
    % \label{fig:}
\end{figure}

\begin{figure}[!ht]
    \centering
    \includegraphics[width=\textwidth]{./images/task4_6.png}
    % \caption{}
    % \label{fig:}
\end{figure}

\begin{figure}[!ht]
    \centering
    \includegraphics[width=\textwidth]{./images/task4_7.png}
    % \caption{}
    % \label{fig:}
\end{figure}

\begin{figure}[!ht]
    \centering
    \includegraphics[width=\textwidth]{./images/task4_8.png}
    % \caption{}
    % \label{fig:}
\end{figure}

\begin{figure}[!ht]
    \centering
    \includegraphics[width=\textwidth]{./images/task4_9.png}
    % \caption{}
    % \label{fig:}
\end{figure}

\clearpage\section{Частина 5}
\label{sec:task5}

\subsection{Завдання}
\label{subsec:task5_task}

Прочитати сигнали кардіоритмограм здорової та хворої людини, збережені в mat-файлі. \\
Навести графіки, позначити вісі. \\
Значення сигналу записуються в мілісекундах,
кожен відлік відповідає інтервалу між ударами серця. \\
Визначити за допомогою спеціальної функції (створити власну) тривалість записаного сигналу,
враховуючи нерівномірну дискретизацію сигналу. \\
Потрібно проінтерполювати сигнал, щоб визначити його значення
на однорідній сітці вузлів для частоти дискретизації 1 Гц. \\
Зберегти отриманий сигнал для використання в подальшій роботі.

\subsection{Код Программи}
\label{subsec:task5_code}
\inputminted{python}{../src/task5.py}

\subsection{Результати}
\label{subsec:task5_results}

\begin{figure}[!ht]
    \centering
    \includegraphics[width=\textwidth]{./images/task5_0}
    \caption{Кардіоритрограма здорової людини}
    \label{fig:hr_healthy}
\end{figure}

\begin{figure}[!ht]
    \centering
    \includegraphics[width=\textwidth]{./images/task5_1}
    \caption{Кардіоритрограма хворої людини}
    \label{fig:hr_apnea}
\end{figure}

\clearpage\section{Частина 6}
\label{sec:task6}

\subsection{Завдання}
\label{subsec:task6_task}

Побудувати функцію, яка будує графік зміни у часі середньої
спектральної густини потужності в заданому частотному діапазоні та часовому
діапазоні для заданого сигналу. В якості параметрів функції
передавати назву сигналу, моменти часу час t1, t2 та частоти f1, f2, а також інші
необхідні параметри

\subsection{Код программи}
\label{subsec:task6_code}
\inputminted{python}{../src/task6.py}

\subsection{Результати}
\label{subsec:task6_results}

\begin{figure}[!ht]
    \centering
    \includegraphics[width=\textwidth]{./images/task6.png}
    % \caption{}
    % \label{fig:}
\end{figure}

\clearpage\section{Частина 7}
\label{sec:task7}

\subsection{Завдання}
\label{subsec:task7_task}

Сформувати сигнал синусоїди частоти 20 Гц амплітуди 1 В та
тривалості 0,5 с для частоти дискретизації 128 Гц. Додати до цього сигналу
шумову складову амплітуди 10 В з нульовим середнім значенням (функція
randn). Розрахувати амплітудний спектр даного зашумленого сигналу.
Побудувати за допомогою функції plot графік сигналу та за допомогою функції
stem його амплітудний спектр.

Збільшувати тривалість зашумленого сигналу (1, 10, 100, 1000 сек.),
отримуючи для кожного сигналу його амплітудний спектр. Побудувати за
допомогою функції plot графіки сигналів та за допомогою функції stem їх
амплітудних спектрів, зробити висновки щодо впливу
тривалості сигналу на роздільну здатність в частотній області та на якість
визначення наявності синусоїдального коливання в шумі за спектральними
характеристиками.

Збільшувати частоту дискретизації зашумленого сигналу тривалості 0,5 с
(1280, 12800, 128000 Гц), отримуючи для кожного сигналу його амплітудний
спектр. На графік кожного разу виводити спектр у межах від 0 до 100 Гц.
Побудувати за допомогою функції plot графіки сигналів та за допомогою
функції stem їх амплітудні спектри, зробити висновки щодо
впливу частоти дискретизації сигналу на роздільну здатність у частотній
області та на якість визначення наявності синусоїдального коливання у шумі за
спектральними характеристиками

\subsection{Код программи}
\label{subsec:task7_code}
\inputminted{python}{../src/task7.py}

\subsection{Результати}
\label{subsec:task7_results}

\begin{enumerate}
    \item Шум замаскував корисний сигнал
    \item Збільшення тривалості запису збільшило роздільну здатність,
          очікуємий пік $20 Hz$ добре розрізняється
    \item Зі збільшенням частоти дискретизації збільшується роздільна здатність сигналу
\end{enumerate}

\begin{figure}[!ht]
    \centering
    \includegraphics[width=\textwidth]{./images/task7_0.png}
    % \caption{}
    % \label{fig:}
\end{figure}

\begin{figure}[!ht]
    \centering
    \includegraphics[width=\textwidth]{./images/task7_1.png}
    % \caption{}
    % \label{fig:}
\end{figure}

\begin{figure}[!ht]
    \centering
    \includegraphics[width=\textwidth]{./images/task7_2.png}
    % \caption{}
    % \label{fig:}
\end{figure}

\begin{figure}[!ht]
    \centering
    \includegraphics[width=\textwidth]{./images/task7_3.png}
    % \caption{}
    % \label{fig:}
\end{figure}

\begin{figure}[!ht]
    \centering
    \includegraphics[width=\textwidth]{./images/task7_4.png}
    % \caption{}
    % \label{fig:}
\end{figure}

\begin{figure}[!ht]
    \centering
    \includegraphics[width=\textwidth]{./images/task7_5.png}
    % \caption{}
    % \label{fig:}
\end{figure}

\begin{figure}[!ht]
    \centering
    \includegraphics[width=\textwidth]{./images/task7_6.png}
    % \caption{}
    % \label{fig:}
\end{figure}

\begin{figure}[!ht]
    \centering
    \includegraphics[width=\textwidth]{./images/task7_7.png}
    % \caption{}
    % \label{fig:}
\end{figure}

\clearpage\section{Частина 8}
\label{sec:task8}

\subsection{Завдання}
\label{subsec:task8_task}

Прочитати сигнал внутрішньо черепного тиску в робочій зоні. \\
Зберегти отриманий сигнал для використання в подальшій роботі. \\
Файл TBI\_ICP.txt - одноканальний сигнал, записаний з частотою дискретизації 125 Гц,
одиниці виміру - мм рт.ст. \\
Вивести графік сигналу на екран, позначити вісі.

\subsection{Код Программи}
\label{subsec:task8_code}
\inputminted{python}{../src/task8.py}

\subsection{Результати}
\label{subsec:task8_results}

\begin{figure}[!ht]
    \centering
    \includegraphics[width=\textwidth]{./images/task8}
    \caption{Внутрішньочерепний тиск}
    \label{fig:intracranial_pressure}
\end{figure}

\clearpage\section{Завдання 9}
\label{sec:task9}

\subsection{Код Программи}
\label{subsec:task9_code}
\inputminted{python}{../src/task9.py}

\subsection{Результати}
\label{subsec:task9_results}

\clearpage\section{Завдання 10}
\label{sec:task10}

\subsection{Код Программи}
\label{subsec:task10_code}
\inputminted{python}{../src/task10.py}

\subsection{Результати}
\label{subsec:task10_results}

\clearpage\section{Висновки}
\label{sec:summary}

\begin{enumerate}
    \item Сформулюйте задачу апроксимації (наближення) функції. \\
          Функцію $f(x)$ необхідно апроксимувати функцією $F(x)$ так,
          щоб відхилення $F(x)$ від $f(x)$ у заданій області було мінімальним
          $|f(x) - F(x)| < \epsilon$;
    \item Яке наближення називають точковою апроксимацією? \\
          Апроксимацію називають точковою, якщо вона побудована на заданій дискретній точок
          $\{x_i\}$;
    \item Яке наближення називають неперервною апроксимацією? \\
          Апроксимацію називають неперервною, якщо вона побудована на неперервній множині точок
          $[a; b)$;
    \item Яке наближення називають інтегральною апроксимацією? \\
          Інтегральна - синонім до неперервна;
    \item Сформулюйте задачу інтерполяції функції. \\
          Побудувати функцію $F(x)$, яка у точках $x_i$ приймає ті самі значення,
          що й задана функція $f(x)$
    \item Наведіть геометричну інтерпретацію задачі інтерполяції. \\
          Побудувати криву $y = F(x)$ певного типу, що проходить через задану дискретну
          множину точок $M_i(x_i, y_i)$;
    \item Що таке глобальна інтерполяція? \\
          Інтерполяційний багаточлен для інтерполяції функції $f(x)$
          побудований на всьому інтервалі зміни аргумента $x$;
    \item Що таке локальна інтерполяція? \\
          Інтерполяційний багаточлен для інтерполяції функції $f(x)$
          побудований окремо для різних частин інтервалу зміни аргумента $x$;
    \item Яке наближення називають екстраполяцією? \\
          Використяння інтерполяційного багаточлена для обчислення значень поза відрізком
          $[x_0, x_n]$;
    \item Які точки називають вузлами інтерполяції? \\
          Точки $x_i, i = \overline{0, n}$ на яких побудована функція $F(x)$
          сприймає ті ж самі значення, що й $f(x)$
    \item Який багаточлен називають інтерполяційним багаточленом Лагранжа? \\
          Многочлен мінімального степеня, що приймає дані значення у даному наборі точок.
          Для $n + 1$ пар $(x_n, y_n)$ існує єдиний многочлен $L(x)$ степеня $n$,
          де $L(x_i) = y_i$
    \item Які багаточлени називають лагранжевими коефіцієнтами?
          \begin{align}
              l_i(x) = \frac{
                  (x - x_0) \dots (x - x_{i - 1})(x - x_{i + 1}) \dots (x - x_n)
              }{
                  (x_i - x_0) \dots (x_i - x_{i - 1})(x_i - x_{i + 1}) \dots (x_i - x_n)
              }
          \end{align}
    \item Наведіть формулу інтерполяційного багаточлена Лагранжа.
          \begin{align}
              L_n(x) = \sum_{i = 0}^n y_i * l_i(x)
          \end{align}
    \item Наведіть формулу інтерполяційного багаточлена Лагранжа
          для випадку лінійної інтерполяції.
          \begin{align}
              L_1(x)
              = \frac{x - x_1}{x_0 - x_1} * y_0
              + \frac{x - x_0}{x_1 - x_0} * y_1
          \end{align}
    \item Наведіть формулу інтерполяційного багаточлена Лагранжа
          для випадку квадратичної інтерполяції.
          \begin{align}
              L_2(x)
              = \frac{(x - x_1)(x - x_2)}{(x_0 - x_1)(x_0 - x_2)} * y_0
              + \frac{(x - x_0)(x - x_2)}{(x_1 - x_0)(x_1 - x_2)} * y_1
              + \frac{(x - x_0)(x - x_1)}{(x_2 - x_0)(x_2 - x_1)} * y_2
          \end{align}
\end{enumerate}

\end{document}
